\documentclass[aspectratio=43]{beamer}
\usetheme{CCNU}

\usepackage{slashed}
\usepackage{datetime}
\usepackage{hyperref}
\usepackage{cancel}
\usepackage[absolute,overlay]{textpos}
\usepackage{tikz}
\usetikzlibrary{arrows.meta,positioning,calc}

\graphicspath{{../latex_docs/figures/}{../utils/screenshots/}}

\def\Put(#1,#2)#3{\leavevmode\makebox(0,0){\put(#1,#2){#3}}}
\def\mydate{\leavevmode\hbox{\the\year-\twodigits\month-\twodigits\day}}
\def\twodigits#1{\ifnum#1<10 0\fi\the#1}
\def\meetingname{Svg2GPU presentation}

\title[Svg2GPU]{Svg2GPU: SVG to GPU rendering}
\author[Andrei Vascul]{Andrei Vascul}
\institute[UBB]{Faculty of Mathematics and Computer Science, Babes-Bolyai University}
\date[\mydate]{\meetingname\\\today}
\titlegraphic{}

\tikzset{
	stage/.style={rounded corners, draw=CCNUMaize!70, fill=CCNUMaize!10, thick, align=center, inner sep=6pt, text width=2.6cm},
	compare/.style={rounded corners, draw=CCNUBlue!50, fill=CCNUBlue!6, thick, align=center, inner sep=6pt, text width=3.35cm},
	arrow/.style={-Latex, very thick, draw=CCNUMaize},
	smallarrow/.style={-Latex, thick, draw=CCNUBlue},
}

\begin{document}

\titlepage

\section{Introduction}

\begin{frame}{Why this exists}
\small
\begin{columns}[T,onlytextwidth]
	\column{0.52\textwidth}
	\begin{itemize}
		\item SVG is a great authoring format: readable, reusable, and semantic.
		\item GPU pipelines want something stricter: triangles, buffers, and draw calls.
		\item \texttt{Svg2GPU} sits in the middle and does the conversion work once.
		\item The goal is practical control, not a full browser-engine rewrite.
	\end{itemize}

	\column{0.46\textwidth}
	\centering
	\begin{tikzpicture}[node distance=1.35cm]
		\node[compare] (svg) {\textbf{SVG}\\paths, groups, styles};
		\node[compare, right=2.25cm of svg] (gpu) {\textbf{GPU}\\buffers, triangles, draws};
		\draw[arrow] (svg) -- node[above, font=\scriptsize]{convert} (gpu);
	\end{tikzpicture}
\end{columns}
\vspace{0.4em}
\begin{block}{In one line}
SVG keeps the intent. GPU code wants the geometry.
\end{block}
\end{frame}

\begin{frame}{SVG path basics}
\small
\begin{columns}[T,onlytextwidth]
	\column{0.58\textwidth}
	\begin{itemize}
		\item Most of the hard work lives in the \texttt{d} attribute.
		\item Commands can be absolute or relative.
		\item \texttt{viewBox} and transforms decide where the shape lands.
		\item Supported primitives stay focused: path, rect, circle, ellipse, line, polyline, polygon, and groups.
	\end{itemize}

	\column{0.40\textwidth}
	\begin{block}{Path commands}
	\centering
	\texttt{M L H V C Q A Z}\\[0.35em]
	\texttt{m l h v c q a z}\\[0.75em]
	\footnotesize Curves and arcs are later flattened into points.
	\end{block}
\end{columns}
\end{frame}

\begin{frame}{WebGPU wants triangles}
\small
\begin{columns}[T,onlytextwidth]
	\column{0.48\textwidth}
	\begin{block}{SVG side}
	\begin{itemize}
		\item declarative scene description
		\item inherited styles and transforms
		\item easy to author, easy to read
	\end{itemize}
	\end{block}

	\column{0.48\textwidth}
	\begin{block}{WebGPU side}
	\begin{itemize}
		\item explicit buffers
		\item vertex formats and indices
		\item shader input and render passes
	\end{itemize}
	\end{block}
\end{columns}

\vspace{0.4em}
\centering
\begin{tikzpicture}[node distance=1.2cm]
	\node[compare, text width=2.95cm] (svg) {SVG says \emph{what} to draw};
	\node[compare, right=1.55cm of svg, text width=3.1cm] (gpu) {WebGPU says \emph{how} to draw it};
	\draw[smallarrow] (svg) -- node[above, font=\scriptsize]{conversion layer} (gpu);
\end{tikzpicture}
\end{frame}

\section{Pipeline}

\begin{frame}{From SVG string to scene}
\small
\begin{center}
\begin{tikzpicture}[node distance=0.95cm]
	\node[stage, text width=2.45cm] (compile) {\texttt{Svg2GPU.compile()}\\entry point};
	\node[stage, right=0.55cm of compile, text width=2.45cm] (parse) {\texttt{SVGParser + StyleResolver}\\typed SVG data};
	\node[stage, right=0.55cm of parse, text width=2.45cm] (build) {\texttt{GeometryBuilder}\\fill / stroke batches};
	\node[stage, right=0.55cm of build, text width=2.35cm] (render) {\texttt{WebGPURenderer}\\upload + draw};

	\draw[arrow] (compile) -- (parse);
	\draw[arrow] (parse) -- (build);
	\draw[arrow] (build) -- (render);
\end{tikzpicture}
\end{center}

\vspace{0.5em}
\begin{block}{Pipeline idea}
Each stage does one job, and the next stage receives already-typed data.
\end{block}
\end{frame}

\begin{frame}{SVGParser}
\small
\begin{columns}[T,onlytextwidth]
	\column{0.55\textwidth}
	\begin{itemize}
		\item walks the XML tree and keeps supported elements
		\item parses presentation attributes and inline \texttt{style}
		\item resolves path \texttt{d} strings into typed commands
		\item rejects missing mandatory data early
	\end{itemize}

	\column{0.43\textwidth}
	\begin{block}{Supported set}
	\centering
	\texttt{path}\\
	\texttt{rect}\\
	\texttt{circle}\\
	\texttt{ellipse}\\
	\texttt{line}\\
	\texttt{polyline}\\
	\texttt{polygon}\\
	\texttt{g}
	\end{block}
\end{columns}
\vspace{0.3em}
\begin{alertblock}{Boundary}
The parser does not render. It only produces typed SVG data.
\end{alertblock}
\end{frame}

\begin{frame}{PathNormalizer}
\small
\begin{columns}[T,onlytextwidth]
	\column{0.52\textwidth}
	\begin{itemize}
		\item relative commands become absolute points
		\item curves are flattened under a tolerance
		\item open and closed contours stay explicit
		\item transforms are applied before tessellation
	\end{itemize}

	\column{0.44\textwidth}
	\begin{block}{Before}
	\centering
	\texttt{M0 0 l10 0 l0 10 z}
	\end{block}

	\begin{block}{After}
	\centering
	\texttt{[(0,0), (10,0), (10,10), (0,0)]}
	\end{block}
\end{columns}
\end{frame}

\begin{frame}{Color and style normalization}
\small
\begin{columns}[T,onlytextwidth]
	\column{0.52\textwidth}
	\begin{itemize}
		\item color input can be HEX, RGB/A, HSL/A, named colors, \texttt{transparent}, or \texttt{none}
		\item WebGPU uses floats from \texttt{0} to \texttt{1}, so conversion happens once
		\item opacity is merged consistently before geometry generation
		\item missing style fields fall back to defaults
	\end{itemize}

	\column{0.46\textwidth}
	\centering
	{\scriptsize
	\begin{tabular}{|l|l|}
	\hline
	\textbf{Input} & \textbf{Normalized RGBA} \\
	\hline
	\texttt{\#00FF33} & \texttt{[0,1,0.2,1]} \\
	\hline
	\texttt{rgba(...)} & \texttt{[1,0,0.5,0.5]} \\
	\hline
	\texttt{hsl(...)} & \texttt{[0,1,0,1]} \\
	\hline
	\texttt{none} & \texttt{[0,0,0,0]} \\
	\hline
	\end{tabular}
	}
\end{columns}
\vspace{0.3em}
\begin{block}{Why this matters}
Boring renderer code is good renderer code.
\end{block}
\end{frame}

\begin{frame}{Tessellation and GeometryBuilder}
\small
\begin{columns}[T,onlytextwidth]
	\column{0.54\textwidth}
	\begin{itemize}
		\item fills become triangles via earcut
		\item strokes become quads with joins and caps
		\item rects, circles, ellipses, polygons, polylines, and lines all funnel into the same geometry path
		\item output is \texttt{GeometryBatch[]} plus stats and diagnostics
	\end{itemize}

	\column{0.42\textwidth}
	\begin{block}{Fill}
	\centering
	contour \(\rightarrow\) triangles
	\end{block}

	\begin{block}{Stroke}
	\centering
	centerline \(\rightarrow\) quads + joins + caps
	\end{block}
\end{columns}
\end{frame}

\section{Rendering}

\begin{frame}{WebGPURenderer}
\small
\begin{columns}[T,onlytextwidth]
	\column{0.53\textwidth}
	\begin{itemize}
		\item creates adapter, device, canvas context, and pipeline
		\item uploads vertex and index buffers per batch
		\item maps the SVG \texttt{viewBox} to clip space using \texttt{fit}
		\item supports resize, MSAA, background color, and cleanup
	\end{itemize}

	\column{0.43\textwidth}
	\centering
	\begin{tikzpicture}[node distance=0.85cm]
		\node[stage, text width=2.2cm] (scene) {scene};
		\node[stage, right=0.55cm of scene, text width=2.3cm] (buffers) {vertex / index\\buffers};
		\node[stage, right=0.55cm of buffers, text width=2.15cm] (pipe) {render\\pass};
		\node[stage, right=0.55cm of pipe, text width=1.95cm] (frame) {frame};
		\draw[arrow] (scene) -- (buffers);
		\draw[arrow] (buffers) -- (pipe);
		\draw[arrow] (pipe) -- (frame);
	\end{tikzpicture}
\end{columns}
\end{frame}

\begin{frame}{Results, demo, and quick start}
\small
\begin{columns}[T,onlytextwidth]
	\column{0.50\textwidth}
	\begin{block}{Quick start}
	\begin{itemize}
		\item \texttt{npm install svg2gpu}
		\item \texttt{Svg2GPU.compile(svg)} for scene output without drawing
		\item docs: \url{http://localhost:5174/docs}
		\item API: \url{http://localhost:5174/typedoc/classes/Svg2GPU.html}
	\end{itemize}
	\end{block}

	\begin{block}{Why it helps}
	It is the fastest route from a repo clone to a usable preview.
	\end{block}

	\column{0.47\textwidth}
	\centering
	\includegraphics[width=\linewidth]{../utils/screenshots/Typedoc.png}
\end{columns}
\end{frame}

\section{Evaluation}

\begin{frame}{Evaluation and limits}
\small
\begin{columns}[T,onlytextwidth]
	\column{0.52\textwidth}
	\centering
	\includegraphics[width=\linewidth]{../utils/screenshots/TigerSVG.png}

	\column{0.44\textwidth}
	\begin{itemize}
		\item simple scenes behave predictably
		\item large SVGs are where the tradeoffs show up
		\item antialiasing is still the hardest part to make perfect
		\item scaling can still favor native SVG
		\item text and gradients are not in the stable path yet
	\end{itemize}
\end{columns}
\vspace{0.3em}
\begin{alertblock}{Takeaway}
The hard parts are visible, which is better than hiding them.
\end{alertblock}
\end{frame}

\begin{frame}{Conclusion and future work}
\small
\begin{itemize}
	\item \texttt{Svg2GPU} makes the SVG \(\rightarrow\) GPU boundary explicit and testable.
	\item The current version is useful where control matters more than perfect SVG parity.
	\item Next steps: broader SVG coverage, stronger contour handling, and cleaner backend abstraction.
	\item Longer term: better rendering quality without losing the current simplicity.
\end{itemize}

\vspace{0.4em}
\centering
\begin{tikzpicture}[node distance=0.85cm]
	\node[compare, text width=2.9cm] (a) {more SVG coverage};
	\node[compare, right=0.9cm of a, text width=2.95cm] (b) {better geometry quality};
	\node[compare, right=0.9cm of b, text width=2.85cm] (c) {cleaner backend split};
	\draw[smallarrow] (a) -- (b);
	\draw[smallarrow] (b) -- (c);
\end{tikzpicture}
\end{frame}

\begin{frame}[plain,t]
\vspace{0.9cm}
\centering
{\Huge Questions?}\\[1.0cm]

\includegraphics[width=0.72\textwidth]{logofooter_csubb}
\end{frame}

\end{document}
